\documentclass[
 aip,
 jcp,
 amsmath,amssymb,
 preprint
]{revtex4-2}

\usepackage{graphicx}
\usepackage{dcolumn}
\usepackage{bm}
\usepackage{hyperref}
\usepackage[T1]{fontenc}
\usepackage[utf8]{inputenc}

\renewcommand{\theequation}{S\arabic{equation}}
\renewcommand{\thefigure}{S\arabic{figure}}
\renewcommand{\thetable}{S\Roman{table}}

\begin{document}

\title{Supplementary Material for:
From Projected Subspaces to Full-Space Implementations:
Representation-Equivalence Audits for Open-Shell ADAPT-VQE}

\author{Lucia Mal\'{\i}\v{c}kov\'a}
\affiliation{Modelos Inteligencia Artificial, Cl. Tajinaste, San Miguel de Abona, Santa Cruz de Tenerife, Spain}

\author{Petr Klenovsk\'y}
\affiliation{Department of Condensed Matter Physics, Faculty of Science, Masaryk University, Kotl\'a\v{r}sk\'a 267/2, 611 37 Brno, Czech Republic}
\affiliation{Czech Metrology Institute, Okru\v{z}n\'i 31, 638 00 Brno, Czech Republic}

\date{\today}
\maketitle

\section{Benchmark Hamiltonians}

\subsection{Type-1 Cu active spaces}

The primary benchmark is an idealized oxidized first-coordination-shell type-1 Cu(II) model with two nitrogen donors, a cysteine-model sulfur donor, and a longer axial sulfur donor.
The molecular orbitals were obtained from a restricted open-shell Hartree--Fock calculation in the STO-3G basis.
The primary active space, CAS(15e,9o), was constructed with PySCF \cite{sun2018pyscf} and the atomic-valence active-space (AVAS) procedure \cite{sayfutyarova2017avas}, targeting the Cu $3d$ and cysteine-S $3p$ manifolds.
CAS(21e,12o) from the same model is used as an expanded algorithmic benchmark.

The active-space Hamiltonian is written
\begin{equation}
\hat H
=
E_{\rm core}
+\sum_{pq} h_{pq} a_p^\dagger a_q
+\frac{1}{2}\sum_{pqrs}(pq|rs)
a_p^\dagger a_r^\dagger a_s a_q .
\label{eq:H}
\end{equation}
The exact Cu target-spin energies used throughout are
\begin{align}
E_D^{18}&=-2518.989067369946~E_h,\\
E_Q^{18}&=-2518.995009103114~E_h,\\
E_D^{24}&=-2518.989131438323~E_h,\\
E_Q^{24}&=-2518.995161621996~E_h.
\end{align}
The lower quartet is a property of these deliberately small approximate active-space Hamiltonians and is not interpreted as a prediction for the experimental oxidized T1 site.

\subsubsection{Canonical Cu geometry and Hamiltonian-generation protocol}

The canonical 18q FCIDUMP can be traced back to the archived geometry-to-Hamiltonian driver.
The idealized first-shell geometry (in \AA) is
\begin{center}
\begin{tabular}{crrr}
\toprule
Atom & $x$ & $y$ & $z$\\
\midrule
Cu &  0.000 &  0.000 &  0.000\\
N  &  2.000 &  0.000 &  0.000\\
H  &  2.340 &  0.900 &  0.200\\
H  &  2.340 & -0.620 &  0.740\\
H  &  2.340 & -0.280 & -0.800\\
N  & -1.000 &  1.732 &  0.000\\
H  & -1.280 &  2.200 &  0.800\\
H  & -1.550 &  2.120 & -0.620\\
H  & -0.120 &  2.260 &  0.040\\
S  & -0.600 & -1.100 &  1.750\\
H  & -1.420 & -1.580 &  2.120\\
S  &  0.200 & -0.500 & -2.830\\
H  &  0.920 & -1.180 & -3.200\\
H  & -0.540 & -0.920 & -3.280\\
\bottomrule
\end{tabular}
\end{center}
The model has total charge $+1$ and spin input $2S=1$.
The short Cu--S donor is represented by SH$^-$ and the longer axial donor by SH$_2$; the two nitrogen donors are NH$_3$ fragments.
These fragments are intentionally minimal and are used only to generate a controlled open-shell benchmark.

The archived PySCF driver uses STO-3G, ROHF with level shift 0.2 and at most 300 SCF cycles (with a second-order Newton fallback if needed), followed by AVAS with labels
\texttt{["Cu 3d","9 S 3p"]} and threshold 0.2.
In the zero-based atom indexing of the geometry above, atom 9 is the short Cu--S donor.
For the returned CAS(15e,9o), a spin-constrained CASSCF is then run with
\texttt{max\_cycle\_macro=50}, \texttt{max\_cycle\_micro=20},
energy tolerance $10^{-6}$, gradient tolerance $10^{-4}$,
FCI-solver tolerance $10^{-7}$, and \texttt{fix\_spin\_(ss=0.75)}.
The reduced one- and two-electron integrals are exported directly to FCIDUMP.
The generator script explicitly allows use of the best orbitals if the macro-iteration cap is reached; a separate convergence flag was not serialized with the archived FCIDUMP, so we do not claim a fully converged chemical CASSCF result from this model.
The scientific use of the file is instead as a fixed, exactly diagonalizable active-space Hamiltonian.

The expanded CAS(21e,12o) FCIDUMP is retained from the same model campaign and is exactly validated as a fixed 24-qubit Hamiltonian.
However, the separate 12-orbital selection driver that originally generated this expanded file is not preserved in the repository.
Accordingly, the 24q case is used only as an algorithmic benchmark and is not claimed to be reproducible from the geometry alone.

\subsection{Independent NO radical benchmark}

As an independent molecular test, an NO radical at $R_{\rm NO}=1.1508$~\AA{} was treated at the ROHF/STO-3G level.
Freezing the lowest four canonical spatial orbitals gives CAS(7e,6o), i.e.\ 12 spin orbitals.
The fixed-$M_S=1/2$ determinant space has dimension 300 and the exact doublet subspace dimension is 210.
The lowest exact doublet energy is
\begin{equation}
E_D^{\rm NO}=-127.632713244383~E_h,
\end{equation}
whereas the lowest quartet is higher by 250.051 m$E_h$.
The lowest doublet is a numerically twofold-degenerate manifold; the next doublet lies 296.438 m$E_h$ higher.
Accordingly, ground-manifold weights rather than fidelity to one arbitrary eigenvector are used when appropriate.

\subsection{Auxiliary OH radical stress test}

An OH radical with O at $(0,0,0)$~\AA{} and H at $(0,0,0.970)$~\AA{} was treated at ROHF/STO-3G, freezing the lowest spatial orbital to obtain CAS(7e,5o), i.e.\ 10 qubits.
The fixed-$M_S=1/2$ space has dimension 50 and the exact doublet subspace dimension is 40.
The lowest doublet is again numerically twofold degenerate; the next doublet lies 223.614 m$E_h$ higher.
This system is used only as an auxiliary stress test because the generalized physical ADAPT search discussed below does not reach the 1.6 m$E_h$ benchmark.

\subsection{Homogeneous Hubbard-chain stress tests}

To remove molecule-specific ingredients, we also use a homogeneous open Hubbard chain
\begin{equation}
\hat H
=
-t\sum_{\langle ij\rangle,\sigma}
\left(
a_{i\sigma}^\dagger a_{j\sigma}
+a_{j\sigma}^\dagger a_{i\sigma}
\right)
+
U\sum_i n_{i\uparrow}n_{i\downarrow}.
\end{equation}
The reported stress tests use four sites, five electrons, $M_S=1/2$, $t=1$, and
$U/t=1,2,4,8$.
They are algebraic validation models, not chemical benchmarks.

\section{Exact spin resolution}

Particle number and $M_S$ do not determine total spin.
We construct
\begin{equation}
\hat S_+=\sum_p a_{p\alpha}^\dagger a_{p\beta}.
\end{equation}
Within the fixed-$M_S=1/2$ determinant space, the pure-doublet highest-weight subspace is
\begin{equation}
\mathcal H_D
=
\ker\!\left(\hat S_+\big|_{\mathcal H_{M_S=1/2}}\right).
\end{equation}
An orthonormal basis $U$ for this null space defines
\begin{equation}
P=UU^\dagger,\qquad Q=I-P,
\end{equation}
and the exact target-space Hamiltonian
\begin{equation}
H_D=U^\dagger H_{M_S}U .
\end{equation}
For a normalized state in the $M_S=1/2$ sector,
\begin{equation}
\langle S^2\rangle
=
\|\hat S_+|\psi\rangle\|^2+M_S(M_S+1).
\end{equation}

The explicit matrix $U$ and projector $P=UU^\dagger$ are classical-oracle objects used because the present benchmark sectors are small enough to resolve exactly.
They are not proposed as scalable quantum-circuit primitives.
For spin, the production physical pool instead imposes the generator-level condition $[B,\hat S^2]\simeq0$ and verifies the exact-projector leakage only as a classical benchmark.
For an arbitrary projected subspace lacking a compact symmetry generator, a scalable surrogate for $QAP=0$ is necessarily problem dependent.

\section{Numerical conventions, tolerances, and optimization settings}

The exact highest-weight basis is constructed with \texttt{scipy.linalg.null\_space} using relative cutoff $10^{-12}$; the stored implementation rejects an $S_+$ null-space defect above $10^{-9}$.
Bare determinant-excitation generators are required to satisfy an anti-Hermiticity defect below $10^{-10}$.
For the generalized physical Cu pools, coefficient null spaces use relative cutoff $10^{-11}$; candidates with Frobenius norm or target-space action below $10^{-12}$ are discarded, and the accepted target-space leakage ratio and normalized $S^2$ commutator defect must each be below $2\times10^{-9}$.
The latter defect is defined as
\begin{equation}
\eta_{S^2}(B)=\frac{\|[\hat S^2,B]\|_F}{\|B\|_F}.
\end{equation}
The reported final pools are many orders of magnitude below these acceptance thresholds.
For the representation audits, a bare generator is classified as numerically non-invariant when $\ell_F>10^{-10}$.

The regenerated canonical Cu projected-ADAPT negative control uses the Hartree--Fock-origin rank-1/rank-2/rank-3 pool without repeated selection.
At each iteration the unused operator with the largest absolute energy derivative is chosen; ties follow the deterministic first-index behavior of \texttt{numpy.argmax}.
All selected amplitudes are globally reoptimized after every addition with L-BFGS-B using \texttt{gtol}$=10^{-8}$, \texttt{ftol}$=10^{-12}$, at most 1000 iterations, and at most 50 line-search steps.
The run terminates when the projected target-doublet error first reaches 1.6 m$E_h$, which occurs at 121 applications.

For the Cu generalized physical ADAPT searches, repeated selection is allowed.
The search implementation reoptimizes the most recent 12 amplitudes after every addition with L-BFGS-B (maximum 80 iterations, \texttt{ftol}$=10^{-12}$, \texttt{gtol}$=10^{-8}$, maximum 30 line-search steps) and performs a global reoptimization every 20 additions (maximum 160 iterations, \texttt{ftol}$=10^{-13}$, \texttt{gtol}$=10^{-8}$, maximum 40 line-search steps), followed by a final global optimization with at least 300 iterations allowed.
Magnitude pruning retains the requested number of applications with the largest $|\theta_k|$ and then globally reoptimizes with L-BFGS-B using \texttt{ftol}$=10^{-14}$, \texttt{gtol}$=10^{-9}$, at most 80 iterations, 50 line-search steps, and memory parameter \texttt{maxcor}$=30$.
These settings describe the stored Cu search and pruning utilities; the reported sequence lengths are therefore search-history dependent rather than proofs of minimality.

All Hartree--Fock tangent ranks reported below are computed by numerical matrix rank with absolute tolerance $10^{-10}$.
Because the FCIDUMP integrals and determinant-space excitation matrices used in these audits are real, the relevant normalized target-state manifold is treated as real: for target dimension $d_D$, the maximum nontrivial tangent rank at one normalized state is therefore $d_D-1$.

\section{Representation-equivalence criterion}

Let $A^\dagger=-A$.
The following conditions are equivalent:
\begin{align}
QAP&=0,\\
[A,P]&=0,\\
e^{\theta A}{\rm Ran}(P)&\subseteq{\rm Ran}(P)
\quad\text{for all real }\theta,\\
P e^{\theta A}P
&=
Ue^{\theta(U^\dagger A U)}U^\dagger
=
e^{\theta PAP}P
\quad\text{for all real }\theta.
\end{align}
If $QAP=0$, anti-Hermiticity also gives
$PAQ=-(QAP)^\dagger=0$, so $A=PAP+QAQ$ is block diagonal.
Conversely, equality of the second derivatives of the last relation at $\theta=0$ gives
\begin{equation}
PA^2P-(PAP)^2
=
PAQAP
=
-(QAP)^\dagger(QAP),
\end{equation}
which vanishes only if $QAP=0$.

The short-amplitude expansions are
\begin{align}
Qe^{\theta A}P
&=
\theta QAP+\mathcal O(\theta^2),\\
Pe^{\theta A}P-e^{\theta PAP}P
&=
-\frac{\theta^2}{2}(QAP)^\dagger(QAP)
+\mathcal O(\theta^3).
\end{align}

\subsection{Finite-amplitude operator-norm bounds}

Define
\begin{equation}
A_0=PAP+QAQ,\qquad V=PAQ+QAP.
\end{equation}
For anti-Hermitian $A$,
\begin{equation}
\|V\|_2=\|QAP\|_2\equiv\lambda.
\end{equation}
The interaction-picture Dyson series gives
\begin{align}
\|Qe^{\theta A}P\|_2
&\le\sinh(|\theta|\lambda),\\
\|Pe^{\theta A}P-e^{\theta PAP}P\|_2
&\le\cosh(|\theta|\lambda)-1.
\end{align}

For a sequence
\begin{equation}
U_{\rm phys}=e^{\theta_MA_M}\cdots e^{\theta_1A_1},
\end{equation}
with block-diagonal surrogate
\begin{equation}
U_0=e^{\theta_MA_{0,M}}\cdots e^{\theta_1A_{0,1}},
\end{equation}
a Duhamel bound and telescoping sum give
\begin{equation}
\|U_{\rm phys}-U_0\|_2
\le
L
\equiv
\sum_{k=1}^{M}|\theta_k|\,\|QA_kP\|_2 .
\label{eq:L}
\end{equation}
Thus $L\ll1$ is sufficient for representation fidelity.
Large $L$ is not a proof of failure because the bound is generally loose.

\subsection{State-specific telescoping certificate}

The operator-norm budget ignores which target-space states are actually visited.
A tighter input-state certificate can be obtained without changing the projected trajectory.
Let
\begin{equation}
V_k=e^{\theta_kA_k},
\qquad
W_k=e^{\theta_kA_{0,k}},
\end{equation}
and let
\begin{equation}
|\phi_{k-1}\rangle
=
W_{k-1}\cdots W_1|\psi_0\rangle
\end{equation}
be the lifted projected prefix.
Define the exact local discrepancy
\begin{equation}
\delta_k
=
\|(V_k-W_k)|\phi_{k-1}\rangle\|
\end{equation}
and
\begin{equation}
D_{\rm seq}=\sum_k\delta_k.
\label{eq:Dseq}
\end{equation}
Using
\begin{equation}
V_M\cdots V_1-W_M\cdots W_1
=
\sum_{k=1}^{M}
V_M\cdots V_{k+1}(V_k-W_k)W_{k-1}\cdots W_1
\end{equation}
and the unitarity of every $V_j$ and $W_j$ gives the rigorous input-state bound
\begin{equation}
\|\Psi_{\rm phys}-\Psi_{\rm proj}\|
\le D_{\rm seq}.
\end{equation}
For normalized states and $D_{\rm seq}\le\sqrt2$,
\begin{equation}
F(\Psi_{\rm phys},\Psi_{\rm proj})
\ge
\left(1-\frac{D_{\rm seq}^2}{2}\right)^2.
\label{eq:Fcert}
\end{equation}
Because $|\Psi_{\rm proj}\rangle\in{\rm Ran}(P)$,
\begin{equation}
p_P(\Psi_{\rm phys})
=
1-\|Q\Psi_{\rm phys}\|^2
\ge
\max(0,1-D_{\rm seq}^2).
\label{eq:Pcert}
\end{equation}
This certificate is state-specific rather than uniform over the complete target subspace.

Table~\ref{tab:statecert} summarizes the molecular negative controls.
The reported diagnostic distance is the phase-minimized pure-state distance
$d_{\rm proj}=\sqrt{2-2|\langle\Psi_{\rm proj}|\Psi_{\rm phys}\rangle|}$.
Because $d_{\rm proj}\le\|\Psi_{\rm phys}-\Psi_{\rm proj}\|$ for any fixed relative phase, it can be compared conservatively with the vector-norm certificate $D_{\rm seq}$ but is not the quantity directly bounded in Eq.~\ref{eq:Dseq}.
\begin{table}[htbp]
\centering
\small
\caption{State-specific sequence certificate. A zero lower bound denotes a mathematically valid but non-informative certificate.}
\label{tab:statecert}
\begin{tabular}{lrrrr}
\toprule
System & $D_{\rm seq}$ & $d_{\rm proj}$ & Certified $F_{\min}$ & Actual $F$\\
\midrule
Cu CAS(15e,9o) & 4.6026 & 1.2075 & 0 & 0.07343\\
NO CAS(7e,6o) & 0.3988 & 0.1678 & 0.8473 & 0.97205\\
OH CAS(7e,5o) & 0.06666 & 0.02776 & 0.99556 & 0.999229\\
\bottomrule
\end{tabular}
\end{table}
For NO, Eq.~\ref{eq:Pcert} also certifies target-doublet weight above 0.8410, compared with the observed 0.97211.
For OH the certified lower bound is 0.99556, compared with the observed 0.999229.
For the severe Cu sequence the state-specific sum is substantially smaller than the uniform operator-norm budget $L=12.686$, but it is still too large to provide a nontrivial fidelity certificate.

\section{Controlled single-generator numerical check}

For a deterministic bare Cu excitation (remove spin orbital 0, add spin orbital 8),
\begin{equation}
\ell_F=
\frac{\|QAU\|_F}{\|AU\|_F}
=
0.5123475383,
\end{equation}
and
\begin{equation}
\|QAP|\Phi_{\rm HF}\rangle\|
=
0.5773502692.
\end{equation}
Exact determinant-space evolutions over $10^{-4}\le\theta\le0.3$ give small-$\theta$ fitted powers
1.0000 for leakage amplitude,
2.0000 for leakage probability, and
2.0000 for the compressed-evolution mismatch norm.
The normalized post-projection infidelity scales approximately as $\theta^{4.03}$ over the numerically resolved part of the same interval.
The resulting scaling is shown in Fig.~\ref{fig:single_scaling}.

\begin{figure}[htbp]
\centering
\includegraphics[width=0.80\linewidth]{figS3_representation_scaling.pdf}
\caption{Single-generator representation-equivalence test in the Cu active space.}
\label{fig:single_scaling}
\end{figure}

As a preserving control, the first generator of the final physical 18q sequence has $\ell_F=3.57\times10^{-16}$ and remains equivalent to double precision for test amplitudes through $\theta=1$.

\section{Canonical Cu projected-ADAPT negative control}

The main negative control was regenerated from the current canonical CAS(15e,9o) FCIDUMP, not inherited from the superseded historical checkpoint.
The deterministic projected ADAPT run contains 121 operators and reaches
\begin{equation}
E_{\rm proj}=-2518.987495463899~E_h,
\end{equation}
or 1.571906 m$E_h$ above the exact target doublet.
Applying the same selected labels and amplitudes to the bare full-space parent generators gives
\begin{align}
E_{\rm bare}&=-2518.524337694489~E_h,\\
\Delta E_{\rm bare}&=464.729675~{\rm m}E_h,\\
\langle S^2\rangle_{\rm bare}&=2.383536856,\\
F_{\rm proj,bare}&=0.0734304134,\\
p_D^{\rm bare}&=0.455487715 .
\end{align}
In this active space only doublet and quartet spin sectors are available in fixed $M_S=1/2$, so
\begin{equation}
p_Q^{\rm bare}=0.544512285 .
\end{equation}

The normalized post-projected doublet component has
\begin{align}
\Delta E_D^{\rm post}&=465.164029~{\rm m}E_h,\\
F(\Psi_{\rm proj},\Psi_{\rm bare,D})&=0.161212720.
\end{align}
The complementary normalized quartet component lies 470.308070 m$E_h$ above the exact quartet minimum.
Thus the failure is not adequately described as contamination of an otherwise correct target state.

The 121 selected generators contain 114 non-invariant factors.
The maximum and median normalized Frobenius leakages are 0.557086 and 0.512348, respectively.
The sequence budgets are
\begin{equation}
L=12.68645,\qquad D_{\rm seq}=4.60259.
\end{equation}
The minimum prefix fidelity is 0.03405 and the minimum target-doublet weight is 0.35760.
The corresponding prefix-level state diagnostics are shown in Fig.~\ref{fig:prefix_state}.

\begin{figure}[htbp]
\centering
\includegraphics[width=0.84\linewidth]{figS4_projected_prefix_state_diagnostics.pdf}
\caption{Prefix-level state diagnostics for the canonical Cu projected-ADAPT negative control.}
\label{fig:prefix_state}
\end{figure}

\section{Fixed parent-sequence full-space reoptimization}

The severe same-angle Cu mismatch tests representation equivalence, but by itself it does not determine whether the 121 parent generators are intrinsically unable to represent a good target-doublet state after their parameters are reoptimized.
We therefore keep the \emph{same 121 full-space generators in exactly the same order} and optimize all amplitudes directly in the fixed-$M_S=1/2$ determinant space.

The target-spin constraint is imposed on
\begin{equation}
p_D(\boldsymbol\theta)
=
\|U^\dagger|\Psi(\boldsymbol\theta)\rangle\|^2
\end{equation}
using SLSQP with analytic derivatives of both the physical energy and $p_D$.
All constrained runs use \texttt{ftol}$=10^{-12}$ and at most 1000 SLSQP iterations; the final constraint value is reported explicitly rather than inferred from the optimizer status.
The determinant-excitation exponentials are evaluated exactly as independent two-level rotations in the determinant basis; this representation agrees with \texttt{expm\_multiply} to better than $1.4\times10^{-16}$ in the validation check.
The operator order is never changed.

For the principal constraint $p_D\ge0.999$, four starts were used: the zero-amplitude vector and three independent Gaussian perturbations of it with standard deviation 0.05 (random seed 9272026).
All four optimizations terminate successfully.
Three converge to the same energy within $5\times10^{-10}$ m$E_h$ and the fourth to a nearby local minimum:
The four principal constrained runs are summarized in Table~\ref{tab:fixed_parent_reopt}.
\begin{table}[htbp]
\centering
\small
\caption{Full-space reoptimization of the fixed 121-generator Cu parent sequence under $p_D\ge0.999$. The projected reduced-space result is 1.5719 m$E_h$.}
\label{tab:fixed_parent_reopt}
\begin{tabular}{lrrrr}
\toprule
Start & $\Delta E$ (m$E_h$) & $p_D$ & $F_{\rm proj}$ & Post-proj.\ $\Delta E_D$ (m$E_h$)\\
\midrule
Zero & 8.481522 & 0.999000 & 0.979887 & 7.335468\\
Zero + perturbation 1 & 8.481522 & 0.999000 & 0.979887 & 7.335467\\
Zero + perturbation 2 & 8.481522 & 0.999000 & 0.979887 & 7.335471\\
Zero + perturbation 3 & 8.973359 & 0.999000 & 0.977998 & 7.995240\\
\bottomrule
\end{tabular}
\end{table}

The best solution has
\begin{equation}
\langle S^2\rangle=0.753000,\qquad
F(\Psi_{\rm proj},\Psi_{\rm full})=0.979887 .
\end{equation}
Tightening the constraint to $p_D\ge0.9999$ and starting from zero amplitudes yields 11.269613 m$E_h$, with $F_{\rm proj}=0.970168$ and a post-projected error of 11.030494 m$E_h$.

The constrained landscape is nonconvex.
Starts in the basin of the original projected amplitudes converge to substantially poorer stationary solutions (43.7706 m$E_h$ for $p_D\ge0.999$ and 47.8917 m$E_h$ for $p_D\ge0.9999$).
These higher minima are retained as an optimizer diagnostic rather than used to support the main conclusion.

This test changes the interpretation of the catastrophic same-angle result.
The fixed parent sequence is capable of approaching the projected target much more closely after target-spin-constrained full-space reoptimization, so the 464.730 m$E_h$ failure should not be described as evidence that the parent sequence is intrinsically useless.
However, none of the converged constrained searches reaches the 1.6 m$E_h$ benchmark, and the best normalized doublet component remains 7.335 m$E_h$ above the exact target.
The robust conclusion is therefore that the reduced-space optimized amplitudes are not transferable to the parent full-space sequence without an equivalence check; a residual fixed-order expressivity/optimization gap remains in the tested searches, but global optimality is not claimed.

To probe fixed-order representability more directly, we also optimize the same 121 amplitudes for fidelity with the lifted projected state rather than for energy.
The analytic fidelity gradient was independently checked against central finite differences at six amplitudes of the canonical checkpoint; the maximum absolute discrepancy was $2.8\times10^{-11}$.
The objective is
\begin{equation}
\max_{\boldsymbol\theta}
F_{\rm proj}(\boldsymbol\theta)
=
\left|
\langle\Psi_{\rm proj}|\Psi_{\rm full}(\boldsymbol\theta)\rangle
\right|^2 .
\end{equation}
Starting from the best $p_D\ge0.999$ energy-constrained basin and three independent perturbations of it, unconstrained fidelity maximization converges in all four cases to
\begin{equation}
F_{\rm proj}=0.988268496,\qquad p_D=0.997155
\end{equation}
to the displayed precision.
The corresponding energy error is 9.8748 m$E_h$ and the normalized doublet component is 7.7646 m$E_h$ above the exact doublet.

Because that fidelity optimum still contains about $0.28\%$ non-doublet weight, we repeat the optimization with the hard constraint $p_D\ge0.999$.
Three tested starts---the best energy-constrained state, a perturbation of it, and the best unconstrained-fidelity state---all converge to
\begin{equation}
F_{\rm proj}=0.987416189,\qquad p_D=0.999000,
\end{equation}
with energy error 10.8166 m$E_h$.
The agreement across these starts identifies a robust local optimum, not a certified global maximum.
Accordingly, the fidelity test strengthens only the empirical statement that the fixed-order full-space sequence does not reproduce the projected state in the searches performed here; it does not prove that exact reproduction is impossible.


\section{Pool-wide invariance and first-order accessibility}

For each bare Hartree--Fock-origin S/D/T excitation, invariance is tested by $\ell_F=0$ to numerical tolerance.
For the Cu CAS(15e,9o) pool, 274 of 323 bare excitations are non-invariant:
22/22 singles, 126/133 doubles, and 126/168 triples.
Only 49 survive simple invariance filtering.

A useful local expressivity diagnostic is the numerical rank of the vectors obtained by acting on the Hartree--Fock reference and restricting to the exact target space.
The rank tolerance is $10^{-10}$.
Because these Hamiltonians, basis vectors, and anti-Hermitian determinant-excitation matrices are real, this audit concerns the real normalized-state manifold; its maximum nontrivial tangent dimension is therefore $d_D-1$ rather than the complex-projective dimension $2d_D-2$.
This is a local first-order diagnostic only and is not a global controllability theorem.
The resulting tangent-space ranks are summarized in Table~\ref{tab:tangent}.

\begin{table}[htbp]
\centering
\small
\caption{Hartree--Fock target-space tangent accessibility. The maximum nontrivial tangent dimension is $d_D-1$.}
\label{tab:tangent}
\begin{tabular}{lrrrr}
\toprule
System & $d_D-1$ & Projected bare S/D/T & Invariant bare & Generalized physical\\
\midrule
Cu CAS(15e,9o) & 239 & 239 & 49 & 99\\
NO CAS(7e,6o) & 209 & 171 & 6 & 56\\
OH CAS(7e,5o) & 39 & 39 & 9 & 25\\
\bottomrule
\end{tabular}
\end{table}

For Cu, simple invariant filtering therefore reduces the first-order accessible target-space rank from the full 239 directions to 49.
The generalized physical rank-1/rank-2 pool restores 99 independent Hartree--Fock tangent directions.
This does not by itself prove global expressivity, but it quantitatively explains why projection can appear highly expressive while simply deleting every leaking bare generator can be too restrictive.

Direct adaptive calculations confirm the point.
Using all 49 invariant Cu bare generators once gives 210.619 m$E_h$ error.
Allowing repeated selection through 100 applications improves only to 167.340 m$E_h$.

\section{Independent NO benchmark}

The NO bare projected S/D/T pool contains 231 generators.
A direct audit finds 225 non-invariant:
17/17 singles, 81/87 doubles, and 127/127 triples.
Projected ADAPT reaches 1.428610 m$E_h$ error in 16 applications.
Eleven of those 16 selected factors are non-invariant.
For the same amplitudes in the bare sequence,
\begin{align}
\Delta E_{\rm bare}&=24.536971~{\rm m}E_h,\\
F_{\rm proj,bare}&=0.972050492,\\
p_D^{\rm bare}&=0.972105986,\\
\langle S^2\rangle_{\rm bare}&=0.833691971.
\end{align}
Here
\begin{equation}
L=0.575942,\qquad D_{\rm seq}=0.398787.
\end{equation}
Unlike the Cu case, post-projection nearly repairs the state:
\begin{align}
\Delta E_D^{\rm post}&=1.477956~{\rm m}E_h,\\
F(\Psi_{\rm proj},\Psi_{\rm bare,D})&=0.999942914.
\end{align}

The lowest exact doublet is twofold degenerate.
The optimized projected state has 0.997999 weight in this manifold, the bare state 0.970221, and the post-projected bare state 0.998061.

Only six of the 231 bare pool elements are invariant.
Using all six once gives 52.154 m$E_h$ error.
Allowing repeats through 60 applications leaves the same energy to the reported precision, showing that simple invariant filtering is again too restrictive.

A generalized physical rank-1/rank-2 pool contains 165 spin-preserving generators.
A 19-application sequence with 13 unique generators reaches
\begin{align}
\Delta E&=1.579026~{\rm m}E_h,\\
\langle S^2\rangle&=0.75,\\
p_D&=1,\\
w_{\rm lowest\ doublet\ manifold}&=0.998967.
\end{align}
The maximum selected representation leakage is below $4.4\times10^{-16}$.

The corresponding 12-qubit Suzuki-2 $r=1$ Qiskit circuit has 7,328 CX gates, depth 9,869, synthesis-energy error 0.002927 m$E_h$, and circuit-to-fermionic fidelity 0.999993871.
It therefore passes the same strict circuit-equivalence criteria used for the Cu circuits.

\section{Auxiliary OH stress test}

For OH CAS(7e,5o), projected ADAPT reaches 1.05508 m$E_h$ in only five factors.
Three are non-invariant, yet the sequence budget is small:
\begin{equation}
L=0.0670781,\qquad D_{\rm seq}=0.0666610.
\end{equation}
The bare same-angle sequence remains close to the projected result:
\begin{align}
\Delta E_{\rm bare}&=1.93276~{\rm m}E_h,\\
F_{\rm proj,bare}&=0.999229332,\\
p_D^{\rm bare}&=0.999229500.
\end{align}
Post-projection yields 1.04156 m$E_h$ error and fidelity 0.999999832 with the projected state.
This provides a useful counterexample to any interpretation that a nonzero invariance defect must always produce a catastrophic error.

The generalized physical pool contains 80 generators and has machine-precision representation leakage.
However, the present first-gradient ADAPT search stalls after repeated reuse of only six unique generators.
After 100 applications,
\begin{align}
\Delta E&=2.58052~{\rm m}E_h,\\
p_D&=1,\\
w_{\rm lowest\ doublet\ manifold}&=0.997527,\\
r&=0.05546~E_h,
\end{align}
while the maximum available first derivative is only $3.9\times10^{-7}$.
The available archived OH closure summary did not preserve the numerical rank tolerance and implementation details required for a reproducible Lie-algebra claim, so no controllability conclusion is drawn here.
The robust conclusion is narrower: a representation-faithful physical pool can still be difficult for this first-gradient adaptive search.
For this reason OH is not counted as a successful physical-ADAPT benchmark in the main text.

\section{Hubbard-model representation stress tests}

The same projected-versus-bare audit was applied to the homogeneous four-site, five-electron open Hubbard chain.
The projected calculation reaches the exact target doublet to numerical precision for every tested $U/t$.
The physical same-angle sequence does not.
The Hubbard results are summarized in Table~\ref{tab:hubbard}.

\begin{table}[htbp]
\centering
\small
\caption{Hubbard-chain stress tests. Energies are in units of $t$.}
\label{tab:hubbard}
\begin{tabular}{rrrrrr}
\toprule
$U/t$ & Ops & Bare error & $F_{\rm proj,bare}$ & $p_D$ & $L$\\
\midrule
1 & 19 & 0.2661 & 0.9244 & 0.9574 & 4.858\\
2 & 19 & 0.5462 & 0.8673 & 0.9656 & 9.149\\
4 & 20 & 1.3914 & 0.6333 & 0.8815 & 9.738\\
8 & 20 & 2.8709 & 0.4763 & 0.9359 & 13.091\\
\bottomrule
\end{tabular}
\end{table}

The corresponding fidelity trends are shown in Fig.~\ref{fig:hubbard}.
\begin{figure}[htbp]
\centering
\includegraphics[width=0.80\linewidth]{figS5_hubbard_representation_audit.pdf}
\caption{Representation-equivalence failure in the homogeneous Hubbard-chain stress tests.}
\label{fig:hubbard}
\end{figure}

These small models are not offered as evidence for chemical relevance.
Their purpose is to show that the representation issue does not rely on the detailed Cu or molecular Hamiltonian.

\section{Full-space spin-preserving operator pools}

This construction is used only as a representation-faithful positive control for the audit.
It is not claimed as a new or preferred scalable spin-adapted ADAPT algorithm; symmetry-preserving state preparation, spin-adapted pools, symmetry-completeness analyses, and exact spin-adapted factorizations are established elsewhere \cite{gard2020symmetry,shkolnikov2023roadblocks,burton2023disco,magoulas2026symmetry,magoulas2025closedform,jain2026factorization}.
Its role here is to provide a full-space sequence for which target-sector preservation can be checked independently and the reduced/full-space trajectories can therefore be compared without representation ambiguity.

A determinant excitation is represented by
\begin{equation}
A=T-T^\dagger.
\end{equation}
Generalized rank-1 and rank-2 $M_S$-preserving excitation patterns are grouped by their spatial remove/add structure.
Within each group, linear combinations
\begin{equation}
B=\sum_j c_jA_j
\end{equation}
are obtained from the null space of the coupling from the target spin space to its complement.
Candidates are retained only when both their target-space leakage and normalized $S^2$ commutator defect lie below numerical tolerance.

For Cu 18q, the complete generalized pool contains 1,080 generators (36 rank-1 and 1,044 rank-2).
The maximum full-pool leakage and normalized $S^2$ commutator defect are at the $10^{-15}$ level.
For Cu 24q, the corresponding pool contains 3,762 generators (66 rank-1 and 3,696 rank-2).
For NO the generalized physical pool contains 165 generators, and for OH 80.

\subsection{Full-sequence positive-control audit}

The final full-space spin-preserving Cu sequences were reconstructed independently in the full fixed-$M_S$ space and in the exact-doublet basis.
For 18q:
\begin{align}
\max_k\ell_F(B_k)&=4.13\times10^{-16},\\
\sum_k|\theta_k|\|QB_kU\|_F&=3.91\times10^{-14},\\
\min_{\rm prefixes}F_{\rm full,reduced}
&=0.9999999999999993.
\end{align}
For 24q:
\begin{align}
\max_k\ell_F(B_k)&=5.17\times10^{-16},\\
\sum_k|\theta_k|\|QB_kU\|_F&=5.14\times10^{-14},\\
\min_{\rm prefixes}F_{\rm full,reduced}
&=0.9999999999999987.
\end{align}
The final full/reduced energies agree to the reported double precision in both cases.
The corrected full-space spin-preserving construction therefore passes the same representation audit that the projected bare sequence fails.

\subsection{Normalization-controlled null-space basis sensitivity}

The spin-preserving coefficient space within one spatial remove/add group can have dimension larger than one, so its individual basis vectors are not physically unique.
The production implementation nevertheless needs a deterministic basis for standard single-operator ADAPT selection and therefore constructs a sparse canonical basis by searching small supports first, fixing a deterministic sign, and snapping coefficients only when they are already within $5\times10^{-10}$ of the exact values $0$, $\pm1$, or $\pm1/\sqrt2$.

A direct normalization audit revealed an important subtlety.
For all 378 two-dimensional Cu 18q null-space groups, the two canonical coefficient vectors have unit Euclidean norm but are not mutually orthogonal:
\begin{equation}
|c_1^\mathsf{T}c_2|=\frac13.
\end{equation}
Likewise, the corresponding determinant-space generators have equal Frobenius norms within numerical precision but normalized Frobenius overlap
\begin{equation}
\frac{|\langle B_1,B_2\rangle_F|}{\|B_1\|_F\|B_2\|_F}=0.3913043478.
\end{equation}
Consequently, applying an $O(2)$ matrix directly to the pair $(B_1,B_2)$ does not preserve the norm of each mixed generator.
Across these groups, such raw mixing can change the Frobenius norm by a factor as large as 1.512 and the Hartree--Fock action norm by a factor as large as 1.732.
Because ADAPT gradients scale with generator normalization, an unnormalized mixing is therefore not a clean test of basis sensitivity.

We instead perform a normalization-controlled prefix audit.
For each of the 51 states along the stored 50-application Cu trajectory (including the Hartree--Fock state), we generate 100 independent random mixes of every two-dimensional null-space group,
\begin{align}
\widetilde B_1&=\cos\alpha\,B_1+\sin\alpha\,B_2,\\
\widetilde B_2&=-\sin\alpha\,B_1+\cos\alpha\,B_2,
\end{align}
and renormalize each mixed direction before comparing the largest ADAPT gradient.
Two conventions are tested independently:
(i) unit Euclidean norm of the underlying coefficient vector and
(ii) the canonical determinant-space Frobenius norm.
All one-dimensional groups are unchanged.

For a metric with Gram matrix $M_{ij}=\langle B_i,B_j\rangle_M$ and signed gradient vector $g_i$, the steepest normalized direction available inside one null-space group has the basis-invariant score
\begin{equation}
G_{\mathcal G}=\sqrt{\boldsymbol g^\mathsf{T}M^{-1}\boldsymbol g}.
\label{eq:group_gradient}
\end{equation}
Under an invertible basis change $B'=BR$, one has $M'=R^\mathsf{T}MR$ and $\boldsymbol g'=R^\mathsf{T}\boldsymbol g$, so Eq.~\ref{eq:group_gradient} is invariant.
The normalization-controlled basis audit is summarized in Table~\ref{tab:basis_sensitivity}.

\begin{table}[htbp]
\centering
\small
\caption{Normalization-controlled basis sensitivity evaluated on all 51 stored Cu 18q prefix states, with 100 independent random bases per prefix. Sensitive prefixes are prefixes for which at least one random basis changes the highest-gradient null-space group relative to the canonical pool.}
\label{tab:basis_sensitivity}
\begin{tabular}{lccc}
\toprule
Metric & Sensitive & Max.\ changed & Group-opt.\ agreement\\
& prefixes & fraction & (prefixes)\\
\midrule
Coefficient Euclidean & 4/51 & 0.14 & 50/51\\
Generator Frobenius & 4/51 & 0.12 & 49/51\\
\bottomrule
\end{tabular}
\end{table}

The controlled result is therefore more modest than a raw basis mixing would imply.
Standard single-operator ADAPT is formally basis dependent, and a few late or near-degenerate decisions can change under normalized rotations, but the highest-gradient group is robust at most sampled prefix states of the stored Cu trajectory.
At the most sensitive coefficient-normalized prefix, 14 of 100 random bases change the top group; under Frobenius normalization the maximum corresponding fraction is 12 of 100.
The canonical single-operator choice agrees with the coefficient-metric invariant group optimum at 50 of 51 prefixes and with the Frobenius-metric optimum at 49 of 51.

These tests do not establish basis-independent full-trajectory convergence because a single altered choice can change all subsequent optimized states.
They do, however, show that the canonical sparse basis does not produce a large normalization-induced selection bias along the stored trajectory.
The reported 50-application passing state remains a result for the documented canonical basis and search/pruning history, not a proof of a basis-independent minimal ansatz.
A group-wise ADAPT rule based on Eq.~\ref{eq:group_gradient} would remove the arbitrary basis choice once a generator normalization metric is specified.

\section{ADAPT selection, repeats, reoptimization, and method comparisons}

For anti-Hermitian candidate $B_k$,
\begin{equation}
g_k
=
\left.\frac{\partial E}{\partial\theta_k}\right|_0
=
\langle\Psi|[\hat H,B_k]|\Psi\rangle
=
2\,{\rm Re}\langle\Psi|\hat H B_k|\Psi\rangle.
\end{equation}
The largest $|g_k|$ is selected, and the pool is not drained, so previously used generators may be selected again, following the original ADAPT construction \cite{grimsley2019adapt}.
Parameters are optimized with L-BFGS-B, including global reoptimization.
Small-amplitude applications are subsequently pruned and the remaining amplitudes reoptimized. This is used only as a practical compaction step and is not presented as a new pruning algorithm \cite{vaquero2025pruned}.

The final Cu 18q sequence has 50 applications and 39 unique generators.
The 24q sequence has 61 applications and 41 unique generators.
The sequence lengths are the smallest passing states found along the present search/pruning histories and are not proofs of global minimality.
The physical-pool comparison is summarized in Table~\ref{tab:poolcomparison}.

\begin{table}[htbp]
\centering
\small
\caption{Cu 18q physical-pool comparison.}
\label{tab:poolcomparison}
\begin{tabular}{lrrrr}
\toprule
Pool & Applications & Error (m$E_h$) & Fidelity & Residual ($E_h$)\\
\midrule
Restricted S+D & 35 & 153.629 & 0.2866 & 0.4189\\
Restricted S+D+T & 147 & 141.392 & 0.2672 & 0.3118\\
Restricted S+D+T + repeats & 250 & 13.166 & 0.8783 & 0.0412\\
Generalized rank-1/rank-2 + repeats & 50 & 1.493 & 0.9826 & 0.0480\\
\bottomrule
\end{tabular}
\end{table}

All rows are spin pure to numerical precision.
This comparison separates spin correctness from practical expressivity.

\section{Qubit mapping and Suzuki synthesis}

Each physical fermionic generator $B_k$ is Jordan--Wigner mapped.
Writing
\begin{equation}
H_k=-iB_k
\end{equation}
gives
\begin{equation}
e^{\theta_kB_k}=e^{-i(-\theta_k)H_k}.
\end{equation}
For
\begin{equation}
H=\sum_{j=1}^{m}H_j,
\end{equation}
the second-order Suzuki step is
\begin{equation}
S_2(t)=
\left(\prod_{j=1}^{m-1}e^{-iH_jt/2}\right)
e^{-iH_mt}
\left(\prod_{j=m-1}^{1}e^{-iH_jt/2}\right).
\end{equation}
With $r$ repetitions, $S_2(t/r)^r$ is used.

The stored validation environment uses Qiskit 2.2.3, Qiskit Aer 0.17.2, and Qiskit Nature 0.7.2.
Circuits are transpiled at optimization level 1 to the abstract basis
$\{R_z,\sqrt X,X,\mathrm{CX}\}$ without a device coupling map.
The stored campaign did not pin \texttt{seed\_transpiler}; the reported CX/depth values should therefore be read as the archived realization for Qiskit 2.2.3 rather than as seed-independent compiler invariants.

The strict circuit-equivalence gate is
\begin{align}
F_{\rm circuit,fermionic}&\ge0.999,\\
|\delta E_{\rm synth}|&\le0.1~{\rm m}E_h,\\
|\langle S^2\rangle-0.75|&\le0.001,\\
p(M_S{\rm\ target})&\ge0.999.
\end{align}

\subsection{Local factorization-error audit}

A finite Suzuki product can violate spin even when the complete fermionic generator commutes with $S^2$, because the constituent Pauli or determinant-excitation terms need not commute with one another or with $S^2$ separately.
We therefore evaluate each ansatz application in isolation at the actual optimized amplitude.

For one Suzuki-2 repetition, the maximum isolated-application infidelity is
\begin{align}
0.1610 &\quad\text{for Cu 18q},\\
2.46\times10^{-4} &\quad\text{for Cu 24q},\\
4.19\times10^{-6} &\quad\text{for NO}.
\end{align}
The 18q sum is dominated by one rank-2 application with $|\theta|\simeq3.13$.

Write one serialized physical generator as $B=\sum_a c_aA_a$, where $A_a$ are its nonzero determinant-excitation constituents (terms with $|c_a|\le10^{-12}$ are ignored for this diagnostic).
In the fixed-$N$, fixed-$M_S$ determinant representation we define
\begin{equation}
c_{\rm comm}(B)
=
\max_{a<b}
\frac{\|[A_a,A_b]\|_F}
{\|A_a\|_F\,\|A_b\|_F},
\label{eq:ccomm}
\end{equation}
with $c_{\rm comm}=0$ for a single active constituent.
This quantity is deliberately simple: it is computed before Jordan--Wigner mapping, ignores the constituent coefficients $c_a$, and is therefore not a rigorous error estimator for the Pauli-level Suzuki formula used by Qiskit.

A direct baseline comparison changes the interpretation of this diagnostic.
For Cu 18q, Cu 24q, and NO, respectively, 24/50, 30/61, and 5/19 applications have $c_{\rm comm}=0$; the largest isolated Suzuki infidelity in those commuting subsets is only $4.4\times10^{-16}$, $1.6\times10^{-15}$, and $2.2\times10^{-16}$.
Thus constituent noncommutativity is an excellent structural separator in these data.
However, inside the subset with $c_{\rm comm}>0$, the Spearman correlations of $|\theta|^3$ alone with isolated-application infidelity are already 0.925, 0.958, and 0.953.
Multiplying by the continuous $c_{\rm comm}$ changes those correlations only to 0.928, 0.970, and 0.953.
We therefore do \emph{not} claim that the numerical magnitude of $c_{\rm comm}$ adds an independent quantitative error law.
The supported conclusion is narrower: commuting constituent sets are easy at this level, while among noncommuting factors the optimized amplitude magnitude carries most of the observed ranking information.
The corresponding noncommuting-factor trend is shown in Fig.~\ref{fig:suzuki_local}.

\begin{figure}[htbp]
\centering
\includegraphics[width=0.80\linewidth]{figS6_local_suzuki_amplitude_v6.pdf}
\caption{Isolated Suzuki-2 infidelity versus optimized amplitude magnitude for applications with $c_{\rm comm}>0$. The amplitude is a strong rank predictor within this noncommuting subset; the continuous $c_{\rm comm}$ factor provides only a modest additional change in the rank correlations reported in the text.}
\label{fig:suzuki_local}
\end{figure}

\subsection{Selective-refinement emulator test}

Because the 18q error is concentrated in a few factors, a determinant-equivalent Suzuki emulator was used to search a restricted nonuniform repetition grid.
The six variables were the applications 9, 19, 21, 25, 33, and 34, chosen from the largest isolated-application errors; every other application was fixed at $r=1$.
The tested repetition sets were
$r_9\in\{1,2,4,8,16\}$,
$r_{25}\in\{1,2,4,8\}$, and
$r_{19},r_{21},r_{33},r_{34}\in\{1,2,4\}$,
for an exhaustive grid of 1620 schedules.
Of these, 277 satisfy the strict state-level criteria.
The lowest component-exponential proxy among this explicitly tested grid refines only applications 9, 21, 25, and 34:
\begin{equation}
r_9=4,\qquad r_{21}=r_{25}=r_{34}=2,
\end{equation}
with all other applications at $r=1$.
The resulting emulator gives
\begin{align}
|\delta E_{\rm synth}|&=0.09502~{\rm m}E_h,\\
F_{\rm circ,fermionic}&=0.9995825,\\
\langle S^2\rangle&=0.750724,\\
p_D&=0.999759.
\end{align}
Its component-exponential proxy is 1600, compared with 1334 for uniform $r=1$ and 5336 for uniform $r=4$.
This result demonstrates concentration of the product-formula difficulty but is \emph{not} reported as a CX/depth reduction:
the selective schedule has not been re-transpiled through Qiskit in the present environment, and transpiler cancellations need not scale linearly with this proxy.

Recent exact and closed-form spin-adapted factorizations \cite{magoulas2025closedform,jain2026factorization} remain important alternatives.
We have not established a one-to-one compiled comparison for the particular numerical linear-combination generators used here.

\section{Measurement grouping and shot allocation}

The Cu 18q and 24q Jordan--Wigner Hamiltonians contain 9,316 and 29,737 Pauli terms.
Greedy QWC grouping gives 1,909 and 6,233 groups.
For
\begin{equation}
\hat H=\sum_g\hat H_g,\qquad
V_g={\rm Var}_\Psi(\hat H_g),
\end{equation}
and independent group sampling,
\begin{equation}
\sigma_E^2=\sum_g\frac{V_g}{n_g}.
\end{equation}
Variance-optimal allocation gives
\begin{equation}
n_g\propto\sqrt{V_g},
\qquad
N_{\rm shot}^{\rm opt}
=
\frac{\left(\sum_g\sqrt{V_g}\right)^2}{\sigma_E^2}.
\end{equation}
This is optimal only for the committed grouping and ideal state preparation \cite{verteletskyi2020measurement}.
The resulting shot estimates are listed in Table~\ref{tab:shots}.

\begin{table}[htbp]
\centering
\small
\caption{Idealized final-energy shot estimates for the chosen greedy QWC grouping.}
\label{tab:shots}
\begin{tabular}{lrrr}
\toprule
& 1.6 m$E_h$ & 1.0 m$E_h$ & 0.5 m$E_h$\\
\midrule
18 qubits & $1.43\times10^8$ & $3.65\times10^8$ & $1.46\times10^9$\\
24 qubits & $5.46\times10^8$ & $1.40\times10^9$ & $5.59\times10^9$\\
\bottomrule
\end{tabular}
\end{table}
At 1.6 m$E_h$, variance-optimal allocation reduces the count relative to equal allocation by factors 26.66 and 42.91 for 18q and 24q, respectively.
These estimates cover final-energy measurement only, not ADAPT gradient rounds or repeated optimizer evaluations \cite{ramoa2025resources}.
Their dependence on target statistical uncertainty is shown in Fig.~\ref{fig:si_measurement}.

\begin{figure}[htbp]
\centering
\includegraphics[width=0.80\linewidth]{figS7_measurement_shots.pdf}
\caption{Idealized final-energy shot estimates for the chosen greedy QWC groupings. The calculation excludes state-preparation error, device noise, calibration overhead, adaptive-gradient measurements, repeated optimizer evaluations, and alternative measurement strategies.}
\label{fig:si_measurement}
\end{figure}

\section{One-particle diagnostics}

The spin-orbital 1-RDM is
\begin{equation}
\gamma_{pq}
=
\langle\Psi|a_p^\dagger a_q|\Psi\rangle.
\end{equation}
For Cu 18q, the spatial 1-RDM Frobenius error is 0.1860, the maximum natural-occupation error 0.002109, and the occupation-vector $L_2$ error 0.002983.
For 24q, the corresponding values are 0.1250, 0.001605, and 0.002985.
The orbital-resolved occupation-number errors are shown in Fig.~\ref{fig:natorb}.

\begin{figure}[htbp]
\centering
\includegraphics[width=0.80\linewidth]{figS1_natural_occupation_errors.pdf}
\caption{Absolute natural-orbital occupation errors of the optimized Cu physical states relative to the exact target doublets.}
\label{fig:natorb}
\end{figure}

The full-state residuals remain 0.0480 $E_h$ and 0.0585 $E_h$, so the compact states should not be interpreted as exact wavefunctions despite their model-space energy accuracy.

\section{Multi-start robustness}

Perturbed-start global reoptimizations finish in a narrow 1.49266--1.49306 m$E_h$ interval for Cu 18q and 1.18070--1.18125 m$E_h$ for Cu 24q.
Several L-BFGS-B runs reach their iteration cap.
These tests therefore support local energetic robustness but do not prove global optimizer convergence or uniqueness.
The individual perturbed-start outcomes are shown in Fig.~\ref{fig:multistart}.

\begin{figure}[htbp]
\centering
\includegraphics[width=0.80\linewidth]{figS2_multistart.pdf}
\caption{Energy errors after reoptimization from perturbed parameter vectors.}
\label{fig:multistart}
\end{figure}

\section{Implementation regression tests}

Five-operator circuit prefixes were used to check fermionic signs, Jordan--Wigner conventions, and circuit assembly.
The Cu 18q and 24q prefixes have circuit-to-fermionic fidelities above 0.999999998 at one Suzuki repetition.

As a negative resource control, the restricted 147-operator Cu S+D+T sequence can itself be synthesized accurately at one Suzuki repetition, but requires 218,502 CX gates at depth 261,514 while its fermionic error remains 141.392 m$E_h$.
A circuit may therefore be a faithful implementation of a poor ansatz.

\section{Excluded Cu--S geometry extension}

The exploratory five-point Cu--S scan failed its pre-specified CASSCF-convergence and active-orbital-continuity gates.
Neither the doublet nor state-specific quartet calculations converged consistently across the full scan, and active-orbital continuity fell below threshold away from the reference geometry.
The scan is excluded from all geometry-dependent scientific claims.

\section{Reproducibility and provenance}

The canonical Cu negative control was freshly regenerated from the current 18q FCIDUMP rather than relying on the missing historical 123-operator checkpoint.
The relevant hashes are listed explicitly below so that the exact files can be identified without relying on filenames alone:

\noindent Cu 18q FCIDUMP:\\
{\small\texttt{4fe2f74860c4a7b1af7e677f62f4d59a1571c2758ff427ba8627607dcd8c8536}}

\noindent 121-operator projected-ADAPT checkpoint:\\
{\small\texttt{d372476940ccb67fd50af104958e36e68f8dddbac40933b9e5b5910e23bd9cd6}}

\noindent 121-operator representation audit:\\
{\small\texttt{f28ee7e78fa790a9c5224bea44b6a5156b1662941be402a1fec712d4facabd92}}

The independent benchmark FCIDUMP hashes used in this revision are:

\noindent NO CAS(7e,6o):\\
{\small\texttt{6c04286a3895f61fef94109baacde752b8ab0c74df0c6642a50919d9f423465b}}

\noindent OH CAS(7e,5o):\\
{\small\texttt{bc9c14bce3d63cda90273cd4c2e3a0ad8138a0be809c635fcb182f7117e5f1e1}}

The public repository is
\url{https://github.com/lucia-malickova/New-VQE-validation}.
The submission-v8 repository state includes the regenerated 121-operator checkpoint and prefix audit, the NO and OH benchmark FCIDUMPs and metadata, the state-specific certificate and tangent-accessibility scripts and outputs, the normalization-controlled null-space-basis audit, the Hubbard stress-test source data, the exhaustive selective-Suzuki search record, the Suzuki diagnostic comparison, and the source data used for the revised figures.
The archived legacy geometry and Hamiltonian-generation scripts are retained for the canonical 18q model; the missing dedicated 24q active-space selection driver is explicitly documented above rather than reconstructed from inference.

\bibliography{references}

\end{document}
